% Reviewer-facing draft for grounding accuracy (Table E).
\subsection{Spatial-Relational Grounding Accuracy}
\label{sec:grounding}

In the deployment loop, STCM is driven by natural-language commands
that must be resolved against the produced graph before Nav2 receives
a goal pose. Sec.~\ref{sec:llm_ablation} isolates the LLM grounder;
this section reports baseline grounding accuracy of the deterministic
template grounder used in the system itself, which is the headline
number a reader should compare against subsequent ablations
(Table~\ref{tab:grounding}). The command sets (\texttt{commands\_meeting.yaml},
$45$ commands; \texttt{commands\_livinglab.yaml}, $30$ commands) are
balanced across four subsets: \emph{simple} (label only),
\emph{disambiguation} (multi-instance, requires choosing among
candidates of the same class), \emph{compositional}
(spatial-relational, e.g.\ ``the trash bin between the doors''), and
\emph{functional} (intent-driven, addressed only via the LLM
ablation in Sec.~\ref{sec:llm_ablation}; not in this table). Each
command is paired across runs (same command IDs, same bag, same
graph-build code, different \texttt{box\_threshold} sweep points
$\in \{0.30, 0.35, 0.40\}$ on the NYU-grounded variant), so the
reported standard deviation reflects threshold-induced graph
variation under deterministic replay rather than seed-level
stochastic replication; we treat $\sigma{=}0\%$ cells as ``replay
repeatability,'' not as ``no run-to-run variance,'' consistent
with the Sec.~\ref{sec:outdoor_caveat} reframing.

\noindent\textbf{Result.}
On the meeting bag, top-$1$ accuracy is $54.0 \pm 5.5\%$ on simple
commands, $34.0 \pm 5.5\%$ on disambiguation, and $34.0 \pm 5.5\%$
on compositional commands, for an overall of $40.7 \pm 5.5\%$
top-$1$ and $42.7 \pm 3.7\%$ top-$2$ over $5$ runs.
On the living lab, top-$1$ falls to $20.0\%$ on simple commands but
recovers on disambiguation ($46.7 \pm 11.5\%$) and compositional
($26.7 \pm 11.5\%$), for an overall of $31.1 \pm 7.7\%$ at top-$1$
and $48.9 \pm 7.7\%$ at top-$2$ -- a substantial top-$2$ recovery
that indicates the template grounder routinely surfaces the correct
candidate in the second slot when an out-of-distribution simple
label trips the first slot. On the outdoor living lab, top-$1$
on simple commands is $20.0\%$, on disambiguation
$26.7 \pm 11.5\%$, and on compositional $20.0\%$, for an overall
top-$1$ of $22.2 \pm 3.8\%$ and top-$2$ of $53.3\%$ on the same
$15$-command spatial subset used for S1 and S2. The S3 perception
ceiling drives the residual gap: a per-command failure-mode
breakdown produced by \texttt{scripts/eval/grounding.py}
attributes $13$/$30$ failures to \texttt{alias\_label\_error} (the
predicted class label does not match the GT canonical label even
when geometry is right, an artefact of the outdoor prompt-bank
vocabulary) and $9$/$30$ to \texttt{perception\_missing} (no
predicted node within the $1.5$~m grounding match radius of the GT
target). Functional / intent commands ($n{=}15$ on S3) are
excluded from the Overall column of Table~\ref{tab:grounding} for
all scenes so that the column denominator is consistent: by
construction the template grounder cannot answer functional
commands (Sec.~\ref{sec:llm_ablation}), so including them on S3
would make S3 incomparable to S1 and S2 (which were scored only on
the same spatial subset). The functional gap on S3 is the same gap
the LLM grounder closes in Sec.~\ref{sec:llm_ablation}.

\noindent\textbf{Interpretation.}
Two failure patterns dominate. First, on living-lab simple commands
the template grounder is bound by the underlying detection ceiling
(Sec.~\ref{sec:semantic}); when the cluttered single-pass trajectory
fails to commit a node for the queried label, no template heuristic
recovers it. This is improved by the NYU-grounded perception
backend (Sec.~\ref{sec:methods_variant}). Second, on disambiguation
the template grounder uses near / between heuristics over object
poses with a small place-graph topology bonus and breaks ties
deterministically; this explains why top-$2 = $ top-$1$ on the
meeting bag (the heuristic does not produce a meaningfully ranked
list for symmetric cases). Sec.~\ref{sec:llm_ablation} shows that
adding an LLM ranker over the same predicted graph closes part of
the disambiguation gap, particularly on functional commands which
the template grounder cannot answer by construction.

% Auto-generated by scripts/eval/render_tables.py
\begin{tabular}{lccccccccc}
  \toprule
  Scene & Runs & Simple top-1 & Simple top-2 & Disambig top-1 & Disambig top-2 & Compose top-1 & Compose top-2 & Overall top-1 & Overall top-2 \\
  \midrule
  meeting & 3 & 36.7 $\pm$ 5.8\% & 76.7 $\pm$ 5.8\% & 36.7 $\pm$ 5.8\% & 50.0 $\pm$ 0.0\% & 40.0 $\pm$ 0.0\% & 60.0 $\pm$ 0.0\% & 37.8 $\pm$ 1.9\% & 62.2 $\pm$ 1.9\% \\
  livinglab & 3 & 60.0 $\pm$ 0.0\% & 80.0 $\pm$ 0.0\% & 33.3 $\pm$ 11.5\% & 33.3 $\pm$ 11.5\% & 53.3 $\pm$ 11.5\% & 53.3 $\pm$ 11.5\% & 48.9 $\pm$ 7.7\% & 55.6 $\pm$ 7.7\% \\
  outdoor\_livinglab & 3 & 20.0 $\pm$ 0.0\% & 60.0 $\pm$ 0.0\% & 20.0 $\pm$ 0.0\% & 20.0 $\pm$ 0.0\% & 60.0 $\pm$ 0.0\% & 60.0 $\pm$ 0.0\% & 33.3 $\pm$ 0.0\% & 46.7 $\pm$ 0.0\% \\
  \bottomrule
\end{tabular}

